\section{Compilación y despliegue}
\label{sec:compilacion}
Para poder compilar y ejecutar la aplicación es necesario disponer de los siguientes requisitos:

\begin{itemize}
    \item Docker CLI o Docker Desktop.
    \item Android Studio.
\end{itemize}

De forma adicional, se recomienda contar con los siguientes elementos para facilitar las tareas de depuración y pruebas:

\begin{itemize}
    \item Python.
    \item Un dispositivo Android físico o una máquina virtual (emulador).
    \item Los puertos \texttt{5432} y \texttt{5050} libres, destinados al backend y al servidor PostgreSQL, respectivamente.
\end{itemize}

\subsection{Pasos para la compilación}

Una vez verificados los requisitos previos, el proceso de compilación y despliegue se realiza siguiendo los pasos que se detallan a continuación:

\begin{enumerate}
    \item Iniciar Docker Desktop o arrancar el \textit{daemon} de Docker.
    
    \item Navegar hasta la carpeta raíz del proyecto y ejecutar el siguiente comando:
    \begin{quote}
        \texttt{docker compose up}
    \end{quote}
    
    \item Obtener la dirección IP de la máquina en la que se va a ejecutar la aplicación\footnote{Esta IP será utilizada por la aplicación móvil para comunicarse con el backend.}.
    
    \item Abrir Android Studio.
    
    \item Ejecutar el proceso de compilación mediante Gradle.
    
    \item Sustituir la IP base configurada en el archivo 
    \begin{quote}
        \texttt{com.nicojero.mysafehaven.remote.RetrofitClient.kt} 
    \end{quote}

    por la dirección IP obtenida previamente\footnote{Es necesario repetir este paso si la IP de la máquina cambia.}.
    
    \item Ejecutar la aplicación en un dispositivo físico o en un emulador Android.
\end{enumerate}
